\section{The Opaque Surface Read}

Once the horizon quotient has been imposed, the grammar must name exactly
what crosses the surface. Too much export would reopen the interior. Too
little export would leave the exterior with no valuation. The chapter uses a
narrow tuple:
\[
  S =
  \bigl(M_{\partial}, Q_{\partial},
        R^{\mathrm{sg}}_{\partial}, \alpha_{\partial}\bigr).
\]
Here \(M_{\partial}\) is the mass surface read, \(Q_{\partial}\) is the
charge surface read, \(R^{\mathrm{sg}}_{\partial}\) is the spin-gravity
residue, and \(\alpha_{\partial}\) is the alpha receipt exported at the
surface.

Each entry has already been earned. Mass is the second-difference strain
read by the loop. Charge is the signed loop count. The spin-gravity residue
is the cost-bearing trace of pushing the surface toward the horizon while
remembering the spin/exclusion drift. Alpha is the finite receipt from the
previous chapter, now carried through the horizon boundary. None of these is
a new interior possession. Each is an exported surface observable.

The tuple is deliberately poor. It does not contain the pairing witness. It
does not contain the spin witness as an inspectable interior fact. It does
not contain the exclusion witness as a hidden mechanism. It does not contain
a complete history of the path by which the interior refined. Those
witnesses may have been meaningful before the horizon. They do not cross
the opaque surface.

The inverse alpha reading may be carried when the alpha receipt permits it,
but it is derived from the same exported receipt. It is not a second path
through the horizon. If \(\alpha_{\partial}\) is exported with enough
finite structure to read an inverse, the inverse belongs to the surface
ledger. If not, the inverse is unavailable for that reading. The grammar
does not invent it to fill a table.

The spin-gravity residue is the most easily misunderstood entry. It is not a
new field hidden in the pair. It is the surface trace of cost: the residue
left when spin-like exclusion and gravity-like horizon approach must be read
together. The residue says that the surface has paid for a correction. It
does not say that the exterior has opened the interior and found the
correction sitting there.

This tuple is the horizon's export contract. Two interiors with the same
\[
  M_{\partial},\quad Q_{\partial},\quad
  R^{\mathrm{sg}}_{\partial},\quad \alpha_{\partial}
\]
are the same exterior read. Two interiors that differ in one of those
surface entries are not the same exterior read. Differences outside the
tuple are masked. Differences inside the tuple are carried.

The holographic illustration is useful if kept grammatical. An admissible
interior cannot claim independent content unless a boundary reconciliation
can carry it. The present chapter uses that discipline narrowly. It does
not claim that every possible interior question has been answered by the
surface. It claims that, for this horizon reading, only the surface tuple is
exported, and valuation must be made from that tuple alone.

The opaque surface therefore converts interior physics into surface
accounting. It permits mass, charge, spin-gravity residue, and alpha to
cross. It forbids the hidden interior witnesses from crossing. It identifies
interiors by equality of the exported tuple. It masks everything else. The
next section gives the name portfolio to the exact surface substitute that
results.
